\mathplain
\section{Related Work}\label{sec:related}

\paragraph{Soundness of DOT.}
The WadlerFest calculus of Amin et al.\ \cite{amin2016essence} is the source of this paper, and its proof characterizes the possible types of a value to invert typing at run time. Rompf and Amin \cite{rompf2016soundness} prove a larger calculus sound by pushing transitivity back through subtyping derivations, with a measure on derivations. Rapoport et al.\ \cite{rapoport2017simple} restrict contexts to inert ones, where every binder has a type of a shape that a value can have. They prove that general typing of a variable can be replaced by tight typing and then by invertible typing, so that the derivations to invert are syntax directed. Amin and Rompf \cite{amin2017definitional} avoid inversion altogether with a semantic model over a definitional interpreter, and Giarrusso et al.\ \cite{giarrusso2020scala} give a step-indexed logical relation in Iris for a DOT with a later modality in its typing rules. Each of these proofs is an argument about derivations of a fixed rule set, and each has been extended by hand when the calculus grew: to paths \cite{rapoport2019path} and to mutation and constructors \cite{kabir2018kdot}. In our proof the derivation is a term, and the only place where derivations are analyzed is the normalizer of Section~\ref{sec:canonical}, which is a function on syntax. Inert contexts correspond to typed stores, which we get from store typing. Invertible typing corresponds to the canonical forms theorem, and the pushback measure to nothing at all, because templates keep composition structural.

\paragraph{Earlier DOT calculi.}
Earlier formulations of DOT \cite{amin2012dot,amin2014foundations} allow application and selection on terms other than variables and have to keep types from mentioning those terms. The WadlerFest calculus moved to a let-normal form, and we keep it, because types that mention variables need names for intermediate results. Direct style is a let-insertion away.

\paragraph{Decidability.}
Subtyping in DOT, and already in its fragment $D_{<:}$, is undecidable \cite{hu2020undecidability}. Algorithmic fragments have been studied for DOT \cite{nieto2017algorithmic} and for related path-dependent systems \cite{mackay2020decidable}. Our checker decides a different problem. It validates fully annotated evidence rather than searching for it, and it is complete for that problem. Producing the evidence is the job of a source type checker, which is necessarily incomplete, and the output of any such checker can be validated independently.

\paragraph{Explicit coercions.}
System FC \cite{sulzmann2007fc} is the model for this work: an intermediate language in which type equalities are proof terms, checked by a small validator that runs after every pass of GHC. Its later versions add roles \cite{breitner2016safe} and generative type abstraction \cite{weirich2011generative}, and its coercions are between types with no term dependence. Coercions as an interpretation of subtyping go back to Breazu-Tannen et al.\ \cite{breazu1991inheritance}, and the coherence of that interpretation for $F_{\le}$ was proven by Curien and Ghelli \cite{curien1992coherence} by rewriting derivations. Our coherence theorem is immediate, because the coercions erase and the erasure is the source term.

\paragraph{Modules and translucent types.}
Type members with bounds are the translucent sums of Harper and Lillibridge \cite{harper1994translucent} and the manifest types of Leroy \cite{leroy1994manifest}, and the F-ing modules translation \cite{rossberg2014fing} interprets ML modules by existential packages in System F. Section~\ref{sec:counterexample} shows where the plain version of that route stops at DOT: a lower bound obtained through another variable's member has to be attached to the original binder, and a package opened at a fresh, unconstrained name does not record it. Whether the sharing machinery of F-ing modules can be adapted to lower bounds we have not investigated. Block names are the replacement for opened names.

\paragraph{Bad bounds elsewhere.}
Stucki and Giarrusso \cite{stucki2021intervals} study higher-order subtyping with type intervals and bad bounds, and their treatment of intervals as pairs of propositions about a name is close in spirit to telescopes over a block. The unsoundness of Scala's and Java's type systems found by Amin and Tate \cite{amin2016unsound} lives in features outside DOT, null and lazy values among them, and our calculus does not speak to it.
